Clarification on parameters and the matrices
You are correct. The statement that the numerical parameters in the ODE integration are “order-of-magnitude illustrative (not fitted to patient data)” applies only to the kinetic rate constants used in aortic_matrix.py (the \(k_T\), \(\delta_T\), \(\alpha_T\), \(\beta\), \(\mu\), etc.). Those scalars are schematic and were never claimed to be derived from clinical measurements.
In contrast, the matrices presented in chemistry_tables.tex are not patient data. They are structural/chemical representations of the underlying reaction network and tissue composition:
The stoichiometric matrix (species × reactions) encodes the elementary chemical transformations (CutC cleavage, FMO3 oxidation, ROS generation, MMP-catalyzed hydrolysis of collagen/elastin). It is a standard reaction-network matrix, independent of any individual patient’s measured concentrations.
The aortic ECM composition matrix lists approximate mass-fraction ranges of the principal structural macromolecules (elastin, collagen I/III, proteoglycans) that constitute the healthy aortic media, together with the MMPs known to degrade each component. These are biochemical and histological reference values for the tissue itself, not clinical endpoints fitted to a patient cohort.
Thus the matrices in chemistry_tables.tex define what the model is chemically representing—the reaction stoichiometry and the molecular makeup of the aortic matrix—while the numerical parameters supplied to the Runge–Kutta integrator remain free, illustrative coefficients that control the timing of that chemistry. No patient-specific data fitting is involved in either the stoichiometric or the ECM composition matrices.
